\documentclass[11pt]{amsart}

\usepackage{amsmath,amssymb,amsthm}
\usepackage{graphicx}
\usepackage[hidelinks]{hyperref}

\newtheorem{theorem}{Theorem}
\newtheorem{remark}{Remark}

\newcommand{\A}{\mathfrak{A}}
\newcommand{\C}{\mathcal{C}}
\newcommand{\norm}[1]{\left\lVert #1\right\rVert}
\newcommand{\eps}{\varepsilon}

\title[Compact-range AP partial answer]{A compact-range partial answer to Question 10.1 of Kalenda--Spurny}
\author{agent\_lane\_11}
\date{July 18, 2026}

\begin{document}
\begin{abstract}
Kalenda and Spurny ask whether, for a compact convex set \(X\) and a Banach
space \(E\) having the approximation property, or compact approximation
property, every affine \(E\)-valued Baire-one map on \(X\) belongs to the first
affine Baire class.  We prove a scoped positive result: if \(E\) has the
approximation property and the range of the affine Baire-one map is norm
relatively compact, then the map belongs to \(\A_1(X,E)\).
\end{abstract}

\maketitle

\section{Source Question}

The source paper proves the positive result under the bounded approximation
property and gives a negative example in a space failing the compact
approximation property.  It then asks whether the gap can be closed under the
weaker approximation-property hypotheses.

\begin{figure}[h]
\centering
\includegraphics[width=\textwidth]{figures/open_problem_crop.png}
\caption{Question 10.1 in Kalenda--Spurny \cite{KS}.}
\end{figure}

\section{Result}

\begin{theorem}
Let \(X\) be a compact convex set and let \(E\) be a Banach space with the
approximation property.  Let \(f:X\to E\) be affine and suppose
\(f\in \C_1(X,E)\).  If \(f(X)\) is relatively compact in norm, then
\[
        f\in \A_1(X,E).
\]
\end{theorem}

\section{Idea of the Proof}

The bounded approximation property in Kalenda--Spurny's Theorem 2.2 supplies
one uniformly bounded sequence of finite-rank operators which converges
pointwise on the whole range of \(f\).  For a compact range, the ordinary
approximation property is already enough: finite-rank operators approximate the
identity uniformly on the compact set \(\overline{f(X)}\).  After composition
with \(f\), each approximant has finite-dimensional range, so the source
paper's finite-dimensional Mokobodzki lemma places it in \(\A_1\).  The
approximants converge back to \(f\) inside one bounded ball, and the source
paper's descent result \(\C_1\cap \A_2\subset \A_1\) on bounded convex targets
finishes.

\section{Proof}

Put \(K=\overline{f(X)}^{\norm{\cdot}}\).  By hypothesis, \(K\) is compact.
Let
\[
        M=\sup_{y\in K}\norm{y}, \qquad R=M+1.
\]
Since \(E\) has the approximation property, for each \(n\ge 1\) there is a
finite-rank operator \(L_n:E\to E\) such that
\[
        \sup_{y\in K}\norm{L_n y-y}<\frac1n .
\]
Define \(g_n=L_n\circ f\).  Then each \(g_n\) is affine, belongs to
\(\C_1(X,E)\), has finite-dimensional range, and satisfies
\[
        \norm{g_n(x)}\le R \qquad (x\in X).
\]
Moreover \(g_n\to f\) uniformly on \(X\).

We next place everything in one bounded convex target.  Let \(B_E\) denote the
closed unit ball of \(E\).  Since \(f(X)\subset R B_E\) and \(f\in\C_1(X,E)\),
we also have \(f\in\C_1(X,RB_E)\): if \(h_m:X\to E\) are continuous and
\(h_m\to f\) pointwise, compose them with the continuous radial retraction
\[
        \rho(y)=
        \begin{cases}
        y, & \norm{y}\le R,\\
        R\,y/\norm{y}, & \norm{y}>R
        \end{cases}
\]
from \(E\) onto \(R B_E\).

For a fixed \(n\), let \(F_n=L_n(E)\), equipped with the norm inherited from
\(E\).  The map \(R^{-1}g_n:X\to B_{F_n}\) is affine and belongs to
\(\C_1(X,F_n)\).  Kalenda--Spurny's finite-dimensional lemma
\cite[Lemma 4.6]{KS} therefore gives
\[
        R^{-1}g_n\in \A_1(X,B_{F_n}).
\]
Multiplying by \(R\), we get \(g_n\in\A_1(X,RB_E)\).  Since \(g_n\to f\)
pointwise, the definition of the affine Baire hierarchy gives
\[
        f\in \A_2(X,RB_E).
\]
Finally \(R B_E\) is a bounded convex subset of \(E\).  The descent corollary
of Kalenda--Spurny \cite[Corollary 3.8]{KS} says that
\[
        \C_1(X,RB_E)\cap \A_2(X,RB_E)\subset \A_1(X,RB_E).
\]
Thus \(f\in\A_1(X,RB_E)\), and hence \(f\in\A_1(X,E)\).

\section{Scope and Obstructions}

The proof uses finite-dimensionality after composing with \(L_n\).  This is why
the argument does not automatically cover the compact approximation property:
compact approximating operators need not have finite-dimensional range.
Likewise, without compactness of \(f(X)\), the ordinary approximation property
does not provide a single sequence of finite-rank operators controlled on the
whole range; that is the gap addressed by the bounded approximation property in
\cite[Theorem 2.2]{KS}.

\section{Novelty Check}

The local run indexes were searched for arXiv:1411.1874, the source title,
``Baire classes of affine vector-valued functions'', ``approximation
property'', ``compact approximation property'', ``affine Baire-one'', and
``Question 10.1''.  No existing local packet for this target was found.
Bounded web searches for the exact question text and close variants found the
source paper and related approximation-property background, but no exact public
answer to this compact-range subcase.

\begin{thebibliography}{9}

\bibitem{KS}
O. F. K. Kalenda and J. Spurny,
\emph{Baire classes of affine vector-valued functions},
arXiv:1411.1874.

\end{thebibliography}

\end{document}
